An archival channeling study compares two collections through an additional computational layer: both are queried using the same UFODex system, and embeddings compare its generated answers to ten matched questions \citep{wahbeh2026records}. These AI-mediated responses differ from the published titles and available captions analyzed here. The cosine values cannot be pooled, and resemblance in either design does not establish an independent source of information.

Tube2Vec is a close computational precedent: it encodes titles and descriptions with all-MiniLM-L6-v2, sums their vectors within each video, and averages the 50 most recent videos per account \citep{boesinger2024tube2vec}. It compares content, social-sharing and recommendation representations using task-specific human judgments. Its filtered sample emphasizes large, mostly English accounts. Our contribution concerns matched text layers and sensitivity in a small, selected corpus; averaging account vectors is not itself novel, and Tube2Vec's human validation does not transfer to ours.

Ethnographic accounts place channeling and mediumship in learned practices and audience interactions \citep{klinoron2014channel,tomlinson2019medium}; our access to these two sources was limited to publisher abstracts. A full-text small-community comparator combines twelve Buddhist practitioners' interviews with reflexive thematic analysis and first-author autoethnography \citep{claisse2023faith}. These different settings motivate contextual interpretation, but our public-record analysis lacks participant accounts or equivalent ethnographic grounding. None supplies a classification of our indexed practitioners.

Topic-evaluation research further distinguishes statistical coherence, useful interpretation and reliability \citep{doogan2021twaddle,hoyle2021incoherence,schroeder2025reliability}. These constructs require different evidence: interpretable labels need examination against documents, while consistency across fitted-model replications does not establish substantive validity. Our rank correlations describe specific perturbations; they are not the replication-reliability coefficients developed for topic models. Human construct validation remains unfinished.
